\documentclass[12pt,a4paper]{article}
\usepackage{ugmath}
\usepackage{float}
\usepackage{placeins}
\usepackage{array}
\usepackage[labelfont=bf,labelsep=space]{caption}
\newcommand{\paren}[1]{\left( #1 \right)}
\newcommand{\alumno}{Ricardo León Martínez}
\newcommand{\materia}{Algebra Lineal I}
\newcommand{\profesor}{Claudia Reynoso Alcántara}
\newcommand{\tarea}{Tarea 5}
\newcommand{\fecha}{2/3/2026}

\begin{document}

    \begin{center}
        {\large\textbf{UNIVERSIDAD DE GUANAJUATO}}\\[0.3cm]
        {\normalsize\textbf{DIVISIÓN DE CIENCIAS NATURALES Y EXACTAS}}\\
        {\normalsize\textbf{CAMPUS GUANAJUATO}}\\[1cm]

        {\Large\textbf{\tarea\ (\materia)}}\\[1cm]
    \end{center}

    \noindent
    \textbf{Nombre:} \alumno 
    \hfill 
    \textbf{Fecha:} \fecha 
    \hfill 
    \textbf{Calificación:} \rule{3cm}{0.4pt}

    \vspace{0.3cm}

    Sea $F$ el campo $\mathbb{Q},\mathbb{R}$ o $\mathbb{C}$.

    \begin{mdframed}[style=mdbluebox,frametitle={Ejercicio 1}]
        Sea $V$ un espacio vectorial sobre $F$. Demuestra que
        \begin{enumerate}
            \item[(a)] El vector cero $0\in V$ es unico.
            \item[(b)] Para todos $v_{1},v_{2}\in V$, existe un único $v_{3}\in V$ tal que
            $v_{1}+v_{3}=v_{2}$.
            \item[(c)] Para todo $v\in V$, se tiene que $(-1)v=-v$ y $0v=0$. 
        \end{enumerate}
    \end{mdframed}

    \begin{mdframed}[style=mdbluebox,frametitle={Ejercicio 2}]
        Considera $V=\left\{(x_{1},x_{2}):x_{1},x_{2}\in\mathbb{C}\right\}$ y define las operaciones:
        \begin{align*}
            +:V\times V&\to V\\
            \left((x_{1},x_{2}),(y_{1},y_{2})\right)&\mapsto(x_{1}+y_{1}+1,x_{2}+y_{2}+1)
        \end{align*}
        \begin{align*}
            \cdot :\mathbb{C}\times V&\to V\\
            \left(c,(x_{1},x_{2})\right)&\mapsto(cx_{1}+c-1,cx_{2}+x-1).
        \end{align*}
        ¿Es $V$ un espacio vectorial sobre $\mathbb{C}$ con estas operaciones? (justifica tu respuesta).
    \end{mdframed}

    \begin{mdframed}[style=mdbluebox,frametitle={Ejercicio 3}]
        Considera el espacio usual $\mathbb{R}^{n}$ sobre $\mathbb{R}$, determina cuáles de los siguientes
        subconjuntos son subespacios vectoriales de $\mathbb{R}^{n}$:
        \begin{enumerate}
            \item[(a)] $W=\left\{(a_{1},\dots,a_{n})\in\mathbb{R}^{n}:a_{1}\geq0\right\}$,
            \item[(b)] $W=\left\{(a_{1},\dots,a_{n})\in\mathbb{R}^{n}:a_{1}+3a_{2}=a_{3}\right\}$,
            \item[(c)] $W=\left\{(a_{1},\dots,a_{n})\in\mathbb{R}^{n}:a_{1}=a_{2}\right\}$,
            \item[(d)] $W=\left\{(a_{1},\dots,a_{n})\in\mathbb{R}^{n}:a_{j}\in\mathbb{Q}\text{ para todo }1\leq j\leq n\right\}$,  
        \end{enumerate}
        justifica tu respuesta.
    \end{mdframed}

    \begin{mdframed}[style=mdbluebox,frametitle={Ejercicio 4}]
        Sean $W_{1},W_{2}$ subespacios vectoriales de un espacio vectorial $V$ sobre el campo
        $F$. Demuestra que si $W_{1}\cup W_{2}$ es un subespacio vectorial $V$ entonces uno de los espacios
        $W_{j}$ está contenido en el otro.
    \end{mdframed}

    \begin{mdframed}[style=mdbluebox,frametitle={Ejercicio 5}]
        Considera el $\mathbb{R}$-espacio vectorial, $V=\left\{f:\mathbb{R}\to\mathbb{R}:f\text{ es una función continua}\right\}$.
        Demuestra que la función $f(x)=\sin(x)$ no es combinación lineal de las funciones $x,x^{2},x^{3}$ (puedes usar todo lo que
        sabes de tus clases de calculo)
    \end{mdframed}

\end{document}